\section{Project Management Plan}
\subsection{Team Work}
\subsubsection{Team Structure and Roles}
The team has no single team leader; responsibility is instead distributed across the five approaches that make up AL-QNN's methodology (Section~4.4), with overlapping support where two approaches are closely coupled. This report is authored from the perspective of the reinforcement-learning module (Approach~4--5).

\begin{table}[htbp]
\centering
\caption{Team structure and roles.}
\label{tab:team}
\begin{tabular}{|p{3.2cm}|p{10cm}|}
\hline
\textbf{Member} & \textbf{Responsibility} \\
\hline
Lâm Vương & Approach~4--5: the RL scheduler (PPO/DQN) and the computational two-simulator design (calibrated surrogate, PennyLane deployment); environment and telemetry integration. \\
Lê Hoàng Nam & Approach~2--3: adaptive circuit scheduling (hardware-efficient ansatz, action space) and safe mutation (soft-masking, commit/rollback registry); barren-plateau measurement (Stage 1B/1C). \\
Đào Quang Minh & Dataset acquisition, preparation, and visualization; supports the Rule and Greedy static baseline schedulers. \\
Trần Bá Thành & Approach~1: multi-dimensional gradient telemetry and diagnosis; supports Lê Hoàng Nam on Approach~2--3. \\
\hline
\end{tabular}
\end{table}

\quad \textbf{Supervisors:} Nguyen Xuan Huy; Nguyen Quoc Trung

\subsubsection{Communication Plan}
The team held weekly progress meetings with the supervisors and communicated daily through a group chat channel. Source code and experimental results were shared through a common repository, and every reported number was required to trace back to a committed CSV log, so that results could be independently re-checked by any team member or advisor.

\begin{itemize}
\item Weekly advisor meeting: progress against the WBS, blockers, and sign-off on any change to the ansatz, reward weights, or evaluation protocol before it was adopted project-wide.
\item Daily internal check-in: short async updates on which stage each member was working on, surfaced early when the RL module (Section~4.7) needed a specific telemetry field from the VQC module (Section~4.4.2) before it could proceed.
\item Shared repository as the single source of truth: every deployment log, comparison table, and figure referenced in this report resolves to a specific committed file (Appendix~D), which is what made the cross-checking practice in the paragraph above enforceable rather than aspirational.
\end{itemize}

\subsubsection{Tools and Repository}
The team used a single shared Git repository for all code, with the quantum-circuit and preprocessing code under the ownership of the VQC-side members and the environment, agent, and deployment code under the ownership of the RL-side member (\Cref{tab:team}). Experimental configuration -- dataset, qubit count, cost type, and backend -- was passed as explicit command-line arguments to the training and comparison scripts listed in Appendix~C, rather than hard-coded, so that a given experiment could be reproduced or varied along a single axis without touching the underlying code. This convention is what made the qubit-count and cost-type sweep reported in Section~6 practical to run as a set of independent, differently-configured invocations of the same two scripts.

\subsection{Project Management Approach}
\subsubsection{Work Breakdown Structure (WBS)}
\begin{table}[htbp]
\centering
\caption{Work breakdown structure (WBS).}
\label{tab:wbs}
\begin{tabular}{|p{2cm}|p{8cm}|p{3.5cm}|}
\hline
\textbf{Phase} & \textbf{Task} & \textbf{Owner} \\
\hline
Stage 1 & Barren-plateau characterization: gradient-variance measurement vs qubits and depth. & Lê Hoàng Nam \\
Stage 1C & Surrogate calibration: fit loss-floor and gradient-variance curves from measured data. & Lâm Vương, Lê Hoàng Nam \\
Stage 2 & RL environment: MDP, telemetry engine, mutation engine, reward. & Lâm Vương \\
Stage 2 & Agent training: PPO and DQN on the surrogate. & Lâm Vương \\
Stage 3 & Deployment: run learned policy on PennyLane; compare against rule and greedy. & Đào Quang Minh \\
Stage 3 & Extended fixed-depth sweep and cost-type regime comparison at the project's 8-qubit scope (Section~6.1--6.2). & Đào Quang Minh, Lâm Vương \\
Rework & Ansatz change (RX-RY-RZ to RZ-RY-RZ) and surrogate re-calibration (Section~2.2.4). & Lê Hoàng Nam, Lâm Vương \\
Rework & Preprocessing standardization: fit-on-train, Z-score before PCA (Section~2.2.4). & Lê Hoàng Nam \\
Support & Datasets, preprocessing pipeline, visualization, demo. & Lê Hoàng Nam, Đào Quang Minh, Lâm Vương \\
Support & Results logging and report documentation. & Trần Bá Thành \\
Exploratory & Telemetry ablation study (with versus without telemetry guardrails, Section~6.1). & Lâm Vương \\
\hline
\end{tabular}
\end{table}

\subsubsection{Risk Management}
\begin{table}[htbp]
\centering
\caption{Risk register.}
\label{tab:risk}
\begin{tabular}{|p{5cm}|p{2.5cm}|p{6.5cm}|}
\hline
\textbf{Risk} & \textbf{Impact} & \textbf{Mitigation} \\
\hline
PennyLane too slow for RL training (millions of steps). & Training infeasible. & Two-simulator design: train on a fast classical surrogate, deploy on PennyLane. \\
Single-seed results are unstable. & Overclaiming. & Report the primary dataset with three seeds; label single-seed results as exploratory. \\
Surrogate diverges from the real circuit after an ansatz change. & Sim-to-real gap. & Re-calibrate the surrogate whenever the ansatz or preprocessing changes. \\
A single fixed-depth sweep on one seed is read as a precise accuracy-vs-depth curve. & Overclaiming a scheduler advantage that is really sampling noise. & Report the sweep as illustrative (Section~6.2) and flag any non-monotonic point explicitly rather than smoothing over it. \\
Data leakage between preprocessing splits. & Inflated accuracy that would not replicate. & Fit every scaler and PCA transform on the training split only, enforced as a single shared preprocessing function (Section~4.2). \\
\hline
\end{tabular}
\end{table}

\subsubsection{Quality Management}
Quality was enforced through three practices. First, data honesty: every figure in this report is intended to trace to a real CSV log produced by the deployment pipeline, and claims that could not be substantiated by logs were removed; one table's original per-seed logs were later found to be unlocatable in the repository, and that gap is disclosed explicitly rather than hidden (Appendix~A.0). Second, a single-source-of-truth code structure: the hardware-efficient ansatz is defined exactly once, so a change there propagates consistently to every stage. Third, multi-seed evaluation on the primary dataset, with explicit caveats wherever only a single seed was available.

\subsubsection{Change Management Process}
Two significant changes occurred during the project and were handled explicitly. The variational ansatz was changed from an RX-RY-RZ rotation block to an RZ-RY-RZ (Euler Z-Y-Z) block; because this preserves three parameters per qubit, parameter counts were unchanged, but the surrogate calibration was scheduled for re-run. The preprocessing pipeline was standardized to fit all scalers on the training split only and to apply a Z-score standardization before PCA, preventing large-scale features from dominating the principal components.

\subsubsection{Closure and Evaluation}
The project closes with a committee defense and this final report. Lessons learned -- particularly the importance of multi-seed evaluation and of honestly reporting cases where the learned agent does not beat a simple baseline -- are captured in the discussion and conclusion.

